\documentclass[12pt, a4paper]{article}

\usepackage[utf8]{inputenc} %Pour l'encodage en utf-8
\usepackage[T1]{fontenc} %Pour les accents
\usepackage[french]{babel} %Pour les langues

\usepackage{siunitx} %Pour les unités
\usepackage{fancyhdr} %Pour les entêtes et bas de page
\usepackage{graphicx} %Pour les figures
\usepackage{mathtools} %Pour les maths
\usepackage[margin=2.5cm]{geometry} %Pour les marges de la page
\usepackage{hyperref} %Pour les liens et les métadonnées

\hypersetup{
	colorlinks=true,
	linkcolor=blue,
	citecolor=blue,
	urlcolor=blue,
	pdfauthor={Maxime BELLAUD, Dalina CHABANE, Khalil FALAH, Paul ROUSSEAU},
	pdftitle={ATCE - Compte rendu}
} %Pour modifier la couleur des liens

\allowdisplaybreaks %Pour avoir une suite d'équations sur plusieurs pages

%Configuration du style de la page
\pagestyle{fancy}
\setlength{\headheight}{14.5pt}
\lhead{ATCE}
\rhead{Compte rendu}

\newcommand{\fc}{\ensuremath{f_c}}

\newcommand{\LT}{\ensuremath{L_T}}

\newcommand{\CT}{\ensuremath{C_T}}
\newcommand{\Cosci}{\ensuremath{C_{\text{osci}}}}

\newcommand{\XLT}{\ensuremath{X_{LT}}}
\newcommand{\XCT}{\ensuremath{X_{CT}}}
\newcommand{\XCosci}{\ensuremath{X_{C \text{osci}}}}

\title{ATCE \\
Asservissement numérique d'un oscillateur à quartz}
\author{Maxime BELLAUD, Dalina CHABANE,\\ Khalil FALAH, Paul ROUSSEAU}
\date{L3 - S6 - 2025 \\
Semaine du 7 au 11 avril 2025}

\begin{document}

\maketitle

{
\hypersetup{hidelinks}
\tableofcontents
}

\section{Introduction}

L'objectif de ce projet est de réguler la fréquence d'un oscillateur à quartz à l'aide d'un microcontrôleur. Autrement dit, un oscillateur à quartz est sensible à son environnement extérieur, notamment la température, et notre objectif est de faire en sorte que la fréquence de l'oscillateur reste stable malgré ces perturbations.

 L'intérêt pour nous est pédagogique car nous réaliserons toutes les étapes pour réaliser ce projet, c'est-à-dire créer le schéma, router le PCB, choisir les composants, souder le circuit, programmer le microcontrôleur, etc. Ce projet nous pousse donc à mettre en œuvre tout ce que nous avons appris pendant nos trois ans de licence, aussi bien en électronique analogique et numérique qu'en automatique.

Ce projet se déroule pendant le dernier semestre de notre licence. Nous réalisons la partie théorique pendant le semestre à l'aide de quelques séances pour nous guider. La réalisation pratique a ensuite lieu pendant une semaine dédiée à ce projet.

Pour réaliser ce projet, nous avons à disposition un PCB routé par nos enseignants ainsi que les composants montés en surface (passifs ou actifs comme le microcontrôleur) pour compléter le circuit. Bien sûr, nous disposons de matériel permettant de les souder. De plus, nous disposons d'oscilloscopes et de multimètres pour réaliser des tests, mais également d'un générateur de signaux afin de simuler un signal 1PPS\footnote{Signal dont le front montant a lieu exactement toutes les secondes} émit par un satellite. Pour terminer sur la partie matérielle, nous pouvons utiliser une imprimante 3D afin de réaliser une boîte pour le PCB.

Les logiciels que nous utiliserons sont pour la plupart gratuits et libres. Pour réaliser le schéma, la simulation de circuits analogiques et le routage du PCB, nous utilisons KiCAD 9.0. Nous utilisons le langage C pour programmer le microcontrôleur et nous utilisons avr-gcc pour compiler les programmes. De plus, le traitement des informations acquises est réalisé avec MATLAB car c'est un logiciel qui nous est familier. Enfin, nous utilisons FreeCAD pour faire de la modélisation mécanique 3D.

Dans un premier temps, nous présenterons comment nous avons conçu le circuit, simulé, puis routé le PCB. Ensuite se trouvera l'estimation du prix des composants et du PCB en fonction des volumes. Après, nous montrerons comment nous avons assemblé le PCB et imprimé par méthode additive une boîte pour celui-ci. Nous aborderons ensuite la programmation du microcontrôleur avant de passer aux tests et à l'évaluation des performances du système.

\section{Schéma du circuit}

% Ceci est un (très) brouillon

Dans un premier temps, nous réalisons le schéma du circuit. L'objectif est de déterminer quels composants choisir et comment les connecter entre eux.

La documentation de l'ATMEGA32U4 peut se trouver à l'adresse suivante : \newline \nopagebreak
\href{https://ww1.microchip.com/downloads/en/devicedoc/atmel-7766-8-bit-avr-atmega16u4-32u4_datasheet.pdf}{https://ww1.microchip.com/downloads/en/devicedoc/atmel-7766-8-bit-avr-atmega16u4- \\ 32u4\_datasheet.pdf}

\subsection{Alimentation}

\subsubsection{Broches}

\begin{itemize}
	\item VCC (alimentation du microcontrôleur) : broches 14, 24, 34 et 44
	\item GND (masse de l'alimentation) : broches 15, 23, 35 et 43
	\item UVCC (alimentation du port USB) : broche 2
	\item UGND (masse du port USB) : broche 5
\end{itemize}

\subsubsection{Découplage}

Ces 5 broches doivent être découplées afin d'avoir une alimentation la plus stable possible. Il nous faut donc calculer la valeur de ces condensateurs.

On sait que $q = C \times U$, donc en dérivant, on obtient $\frac{dq}{dt} = C \times \frac{du}{dt}$. De plus, on sait que $\frac{dq}{dt} = i$ donc $i = C \times \frac{du}{dt}$, que l'on peut aussi écrire $i = C \times \frac{\Delta u}{\Delta t}$ donc $C = \frac{i \Delta t}{\Delta u}$.

Sur la figure 30-4 de la documentation de l'ATMEGA32U4, on lit une consommation de courant du microcontrôleur de \SI{13}{\milli\ampere} pour \SI{5}{\volt} et \SI{16}{\mega\hertz}. De plus, $\Delta t = \frac{1}{\SI{16}{\mega\hertz}} = \SI{62}{\nano\second}$. Enfin, on choisit une fluctuation de tension $\Delta u = \SI{0.1}{\volt}$.

On obtient alors $C = \SI{8.125}{\micro\farad}$. Nous choisissons alors d'arrondir à $\boxed{C = \SI{10}{\micro\farad}}$, ce qui correspond alors à un courant de \SI{16}{\milli\ampere}, toujours pour $\Delta V = \SI{0.1}{\volt}$

\subsection{Bouton de reset}

Le signal pour réinitialisé le microcontrôleur se trouve sur la broche 13. De plus, on voit sur la figure 8-1 de la documentation de l'ATMEGA32U4 qu'une résistance de pull-up est déjà placée dans le microcontrôleur. Il suffit donc de placer un bouton entre la broche 13 et GND.

\subsection{Bouton d'activation du bootloader}

Le signal d'activation du bootloader est noté HWB et est actif à l'état bas. Contrairement à la broche reset, il n'y a pas de circuit de pull-up intégré. On réalise alors un circuit pull-up standard afin que la broche soit à l'état haut lorsque le bouton est relâché et qu'il ne reste pas dans un état flottant. Nous choisissons une résistance de \SI{5}{\kilo\ohm}, qui est une valeur standard pour ces circuits.

\subsection{Connecteur coaxial pour signal GPS}

Le signal GPS générera un signal d'une impulsion par seconde (1PPS). Il faudra mesurer cette période à l'aide d'un Input Capture (IC). Nous choisissons le timer 3, qui est un timer sur 16 bits afin d'avoir le plus grand nombre de bits possibles. Ainsi, l'ICP3 correspond à la broche 32. Nous pensons également à raccorder le blindage du connecteur à la masse.

\subsection{Connecteur USB}

Connecteur de type mini-USB pour pouvoir utiliser le même câble que la Olimexino.

\begin{itemize}
	\item VBUS du microcontrôleur (broche 7) connecté au VBUS du connecteur USB. C'est l'alimentation \SI{5}{\volt}.
	\item D+ et D- du microcontrôleur (broches 3 et 4) connectées au D+ et D- du connecteur USB. C'est une paire différentielle permettant de transmettre des données. Des résistances de \SI{22}{\ohm} sont placées entre les D+ et D- du microcontrôleur et du connecteur USB (partie 2.2.8 et 2.2.9 de la documentation de l'ATMEGA32U4).
	\item Le blindage (shield) est connecté à la masse.
	\item ID du connecteur USB laissé déconnecté pour indiquer que le microcontrôleur est un périphérique (\href{https://www.astuces-pratiques.fr/informatique/usb-otg-broche-id-et-cable}{https://www.astuces-pratiques.fr/informatique/usb-otg-broche-id-et-cable}).
\end{itemize}

Ajout d'une résistance de \SI{0}{\ohm} entre VBUS et le reste du circuit. Cela permettra de s'assurer que l'on a bien \SI{5}{\volt} avant de connecter le reste du circuit.

\subsection{Convertisseur numérique-analogique}

Pour réaliser un DAC, nous utilisons un signal PWM suivi d'un filtre passe-bas.

\subsubsection{PWM}

Utilisation du timer 1 car il a une taille de 16 bits. Cela permettra d'avoir plus de liberté qu'un compteur sur 8 bits. La sortie du PWM sera donc le OC1A (broche 29).

\subsubsection{Filtre passe-bas}

La fréquence de coupure d'un filtre passe-bas est donnée par $\fc = \frac{1}{2 \pi R C}$. En prenant les valeurs des composants les plus utilisés, c'est-à-dire $C = \SI{10}{\micro\farad}$ et $R = \SI{5}{\kilo\ohm}$, on obtient $\fc = \SI{3.18}{\hertz}$. Ainsi, en générant un signal PWM d'au moins \SI{1}{\kilo\hertz}, le signal devrait être filtré correctement et seule la composante continue sera conservée. La fréquence du PWM peut être augmentée pour réduire l'oscillation de la tension en sortie du filtre (voir avec les simulations).

\subsection{Résonateur}

Connexion de l'oscillateur entre XTAL1 et XTAL2 comme indiqué sur la figure 6-2 de la documentation de l'ATMEGA32U4.

\subsubsection{Varicap}

Afin de modifier la fréquence de résonance de l'oscillateur à quartz, on utilise une Varicap pour modifier la valeur des capacités de pied C1 et C2.

\subsubsection{T de polarisation}

La polarisation de la Varicap ne doit pas influencer le résonateur et le signal de l'oscillateur ne doit pas perturber le circuit de polarisation de la Varicap. Pour cela, on utilise un T de polarisation.

La capacité du T de polarisation notée \CT\ ne doit pas influencer la valeur des capacités de pied. Ainsi, on doit avoir $\XCT \ll \XCosci \Leftrightarrow \frac{1}{\CT \omega} \ll \frac{1}{\Cosci \omega} \Leftrightarrow \CT \gg \Cosci$.

\Cosci\ dépendra du résonateur choisi mais est de l'ordre de \SI{20}{\pico\farad	}. Ainsi, on peut choisir $\CT = \SI{1}{\micro\farad}$.

Pour la bobine du T de polarisation, à \SI{16}{\mega\hertz}, il faut que $\XLT \gg \XCT \Leftrightarrow 2 \pi f \LT \gg \frac{1}{2 \pi f \CT} \Leftrightarrow \LT \gg \frac{1}{4 \pi f^2 \CT^2}$, avec $f$ la fréquence en \si{\hertz}. On choisit donc $\LT = \SI{10}{\micro\henry}$.

Pour vérifier, à très basse fréquence, on a bien $\XLT \to 0$ et $\XCT \to \infty$. À \SI{16}{\mega\hertz}, on a $\XLT \approx \SI{1005}{\ohm}$ et $\XCT \approx \SI{10}{\milli\ohm}$.

\subsubsection{Suiveur}

Afin de s'assurer que la tension de polarisation de la diode est imposée par le DAC, on place un amplificateur monté en suiveur entre la sortie du filtre passe-bas et le T de polarisation de la Varicap (du côté de la bobine). Le suiveur est un rail-to-rail afin de pouvoir exploiter la plage de valeur complète de 0 à \SI{5}{\volt}.

\subsection{Autres}

\subsubsection{UCAP}

La documentation du microcontrôleur (partie 2.2.12) indique de connecter UCAP à la masse via un condensateur de \SI{1}{\micro\farad}.

\subsubsection{AREF}

Étant donné que le projet ne nécessite pas de conversion analogique-numérique, l'ADC ne sera pas utilisé et AREF (broche 42) peut être laissée déconnectée.

\subsection{Points de test}

Afin de vérifier que le circuit fonctionne comme prévu, nous plaçons des points de test qui permettront de visualiser la tension à ces points.

La résistance de \SI{0}{\ohm} fait office de point de tests pour l'alimentation. La résistance de \SI{0}{\ohm} sera placée une fois que l'on aura vérifié qu'un potentiel de \SI{5}{\volt} est bien présent à cet endroit.

Nous ajoutons 3 points de test, un à la sortie PWM, un à la sortie du filtre passe-bas, et un à la sortie du suiveur. Ainsi, cela nous permettra de vérifier que le signal PWM est correct, que le signal est ensuite bien filtré, et enfin que l'amplificateur suiveur fonctionne.

Nous choisissons de ne pas mettre de point de test vers le résonateur pour ne pas ajouter de bruit au signal.

\section{Choix des composants}
\subsection{Microcontrôleur et USB}
L’Atmega32U4 a été choisi pour sa capacité à gérer les fonctions d’asservissement nécessaires. Il intègre notamment des timers 16 bits, essentiels pour l’input capture afin de mesurer précisément la fréquence de l’oscillateur, et permet l’émission d’une PWM 10 bits pour le filtrage du signal de commande. \\

Les critères de choix pour l'USB sont d'abord le format (mini ou micro) dépendant des préférences d’encombrement et de robustesse mécanique, puis le connecteur horizontal ou vertical selon la disposition de la carte. L'empreinte “Horizontal” : ce modèle (souvent référence “USB\_OTG\_SMD\_xxx”) se soude facilement et supporte un alignement correct sur la carte.

\subsection{Résonateur à quartz et Varicap}
Pour l'empreinte du quartz à 16MHz, on choisit SMD Abracon ABM7 (généralement 5,0×3,2mm) : très répandu, faible encombrement, montage en surface pour simplifier la fabrication. \\
Il est suffisamment précis pour le microcontrôleur et la mesure du temps, et compatible avec l’asservissement GPS, comme il est assez répondu chez Farnell, Mouser, etc. \\

La BB833 est une diode varicap (ou varactor) utilisée pour ajuster la fréquence d’oscillation du quartz. Etant un ancien modèle pas forcément disponible chez les fournisseurs, on opte pour l'alternative du BB857 qui a des caractéristiques similaires : plage de capacité adaptée à la dérive de fréquence que l’on souhaite corriger, tension maximale (ou tension inverse de fonctionnement) compatible avec l’alimentation choisie, empreinte SMD (boîtier 0805/2012 Metric) : format compact, facile à souder en surface, et disponible chez de nombreux fournisseurs.

\subsection{Boitiers et connecteurs}
Le format 0805 offre une taille réduite, ce qui permet de réaliser une densité élevée sur le PCB tout en conservant une bonne maniabilité lors de l’assemblage et soudage. Pour les composants nécessitant une puissance ou des tolérances particulières, le format 1206 peut être utilisé. Ces boitiers sont standards sur le marché, ce qui facilite leur approvisionnement auprès de fournisseurs reconnus. \\



\subsection{Liste des composants (BOM)}
Pour chacun des éléments (condensateurs, résistances, bobines, diode varicap, quartz, etc.), nous avons consulté les fiches techniques disponibles sur Farnell afin de :
\begin{itemize}
\item Vérifier la valeur électrique (par exemple, 1 $\mu F$ pour un condensateur de découplage, 22 $\Omega$ pour certaines résistances ou les valeurs spécifiques pour le réseau de charge du quartz).
\item Confirmer que les caractéristiques dimensionnelles du composant correspondent à l’empreinte utilisée dans notre design KiCad. Par exemple, le choix entre un boîtier CMS en format 0805 ou 1206 dépend de la densité de montage voulue et des contraintes mécaniques lors de l’assemblage.
\item Chaque composant a été comparé en fonction de son prix unitaire. Par exemple, un condensateur de 1 $\mu F$ en 0805 (référence 2470435) a été retenu car son prix est de l’ordre de 0,118 $\euro $ et il est parfaitement adapté aux besoins de découplage du microcontrôleur.
\item La disponibilité du composant sur Farnell ainsi que la possibilité d’en commander en quantité suffisante ont également été prises en compte.
\end{itemize}


\section{Routage du PCB}

\section{Estimation du coût}

\section{Conclusion}

%\cite{Référence} pour citer la référence. La référence est l'argument de bibitem. Par exemple pour bibitem{Toto}, on fait \cite{Toto}

\section{Bibliographie}
\renewcommand{\refname}{}
\vspace{-1cm}
\begin{thebibliography}{}
	\bibitem{Technique de soudure de composants CMS}
		Technique de soudure de composants CMS : \url{https://youtu.be/us5zLlp1aV0?si=Flcb0uajIgZAhLWc}
\end{thebibliography}

\newpage
\phantom{.}
\vfill
\begin{center}
\section{Annexes}
\end{center}
\vfill
\newpage

\end{document}